\scanpage{62}
\mspage{11}

\deleted{Comme $X$ et $Y$ sont homotopiquement lisses et propres sur $I$, donc des $W$-fibrés sur $I$ par $\textcircled{4}$, on voit que $f_i:X_i\to Y_i$ est dans $W$ pour tt $i\in I$. Donc \unsure{OBS|OPS} $0$-connexes.}

% slanted
\margin{À nouveau, on \unsure{utilise|voit} un \unsure{ingrédient} \unsure{mystérieux}, \unsure{inutile} : \uline{tout \unsure{baigne}}. en fait, \unsure{faut} : il faut $W_\infty$-parfait (i.e. $W_\infty$-lisse et propre) et $W$-parfait [en fait, $W_\infty$-parfait]}

Je factorise donc $f$ en
\[
 X\xrightarrow{i_{\unsure{*|X}}}\bar X\xrightarrow{\bar f}Y,
\]
\deleted{\unsure{cet}} $\bar f$ une $W_\infty$-fibration\inserted{et une $W$-fibration$^*$}, $i_X\in W\cap W_\infty$ --- cela marche fonctoriellement, par la construction de $\overline{\underline{\mathrm{VII}}}$ p.170-171 (du Cartan, Serre), à partir des propriétés de $Y\to\underline{\mathrm{Ch}}(Y)\to Y\times Y$ sans plus.

\deleted{Donc $f\in W\cap W_\infty$} Comme $\bar f i=f\in W$\inserted{et $i\in W$}, on déduit $\bar f\in W$. \deleted{Donc pour prouver $f\in W_\infty$} Comme $i\in W_\infty$, on en \unsure{conclut|déduit} $f\in W_\infty\Longleftrightarrow\bar f\in W_\infty$.

Donc pour prouver que $f$ est dans $W_\infty$, on est ramené au cas où $f$ est une $W$-fibration et une $W_\infty$-fibration. Mais alors la prop. 4 nous montre que les fibres de $f$ sont $W$-asphériques, et il suffit de voir qu'elles sont $W_\infty$-asphériques (prop. $3^*$ appliquée à $W_\infty$),

% slanted
\margin{on n'a utilisé ici que \unsure{cette} propriété, sur les fibres, comme conséq. de $\textcircled{4}$ (ou utiliser 4 bis ou sa forme la plus faible $4\,\mathrm b''$ cf. pages 13, 14.}

\margin{on n'utilise la prop 3 que \unsure{pour|dans} le cas $Y=S$, pour lequel $W(2~\mathrm{bis})$ suffit}

Ainsi on est ramené au cas où $Y=e$, i.e. à prouver que si $X$ dans $\mathrm{Cat}$ est $W$-asphérique, il est $W_\infty$-asphérique. \deleted{L'axi} \deleted{L'axiome de connexité sur $W$ nous assure déjà que $X$ est $0$-connexe.}

Soit $x_0\in\operatorname{Ob}X$, considérons comme \deleted{\unsure{tantôt}}
\[
 E=\underline{\mathrm{Ch}}(X,x_0,-)\xrightarrow{p}X,
\]

on sait que $E$ est $W\cap W_\infty$-asphérique, et que c'est une fibration pour $W$ et pour $W_\infty$. Comme $X$ est $W$-asphérique\inserted{(comme $E$)}, $p\in W$ \deleted{\unsure{est}}, donc \inserted{à nouveau} par la prop. $4^*$ les fibres de $p$ sont $W$-asphériques, en particulier (par $\textcircled{5}$) $0$-connexes. On a donc prouvé le

% slanted
\margin{en fait, la condition $\unsure{W(4\,\mathrm b'')|W4\,\mathrm b''}$ sur $W$}

% archive sheet number: 62
